% Ad Atticum 15.15 (ad-atticum-15-15) --- English translation
% Translated by Claude (Anthropic) from the Latin (Perseus).
% Style guide: STYLE.md.

\ciceroLetterOpener{CICERO ATTICO salutem}

\ciceroSection{1} A curse on Lucius Antonius, if he really is making trouble for the Buthrotians! I have drawn up a deposition which will be sealed whenever you wish. As to the money of the Arpinates, if Lucius Fadius the aedile demands it, hand over the whole sum. I have written to you in another letter about the 110,000 sesterces which were to be looked after by Statius. So if Fadius asks for it, I want it paid to him; to no one but Fadius. I have written to Eros that he should pay it over from a deposit \dag I believe is held with me. \dag

\ciceroSection{2} I detest the queen. That I do so with reason, Ammonius, the guarantor of her promises, knows --- promises which were \textit{literary} \textit{[Greek: philologa]} and worthy of my standing, such that I should have dared to speak of them even in a public assembly. As for Sara, besides being a wicked fellow, I have found him insolent toward me into the bargain. I saw him at my house once and once only. When I asked him \textit{cordially} \textit{[Greek: philophronos]} what he wanted, he said he was looking for Atticus. The arrogance of the queen herself, when she was in her gardens across the Tiber, I cannot recall without great pain. Nothing, then, with those people. They suppose I have no spirit --- or rather no stomach.

\ciceroSection{3} My departure, as I see, is being held up by Eros's bookkeeping. For although by the balance of what he made up on the 5th of April I ought to be in surplus, I am being forced to borrow; and what was received from those productive sources of yours, I had reckoned set aside for that shrine. But I have entrusted this business to Tiro, whom I have sent to Rome on that account; I did not wish to encumber you when you are encumbered already.

\ciceroSection{4} Our Cicero, the more modest he is the more he moves me. For he has written nothing to me about this matter --- to whom of all people he most surely should have written; he has written instead to Tiro that nothing has been paid him since the Kalends of April (for the year's allowance is reckoned to run from then). It has always been your view, in keeping with your nature, and you have judged \dag that it bears upon my dignity \dag that he be treated by us not only most liberally but lavishly and abundantly. I should therefore like you to see to it (and I should not be troubling you, were there any other way I could manage this) that his year's allowance be transferred for him to Athens. Eros will of course count out the money. It is for this purpose that I have sent Tiro. So you will take care of it, and write to me whatever seems right to you concerning him.
